\section{Introduction}
\label{sec:intro}

\begin{figure}[h]
\centering
\includegraphics[width=\textwidth]{signature_mechanism}
\caption{\textbf{Why \ours{} works.}
\textbf{(A)} Training 5 seeds of the same model at $T_\mathrm{max}=100$
and 5 more at $T_\mathrm{max}=150$ reaches nearly identical final
validation accuracy (0.6771 vs.\ 0.6776, $<\!0.001$ apart).
\textbf{(B)} Yet the two endpoints \emph{do not} lie in the same loss
basin: linearly interpolating their weights crosses a $+0.25$~nats
train-loss barrier, and test accuracy dips by $6$~percentage points in
the middle.
\textbf{(C)} On the 13\% of test samples where the 100ep group-consensus
and 150ep group-consensus \emph{disagree}, the mixed ensemble is
\textbf{+10.3~pp} more accurate than either single-length consensus. The
macroscopic $+1.5$~pp ensemble gain is the product of this correction
($10.3\%$) and the disagreement fraction ($13.1\%$).}
\label{fig:hero}
\end{figure}

Deep ensembling \cite{lakshminarayanan2017simple} remains one of the most
effective and widely used uncertainty-estimation and accuracy-boosting
techniques in deep learning. The classic recipe is simple: train $N$
independent models with different random seeds, average their softmax
outputs at test time, and enjoy a $\sim$1--3 percentage point accuracy boost
in exchange for $N\times$ inference cost. But seed diversity is a blunt
instrument, and practitioners have spent a decade looking for richer
ensemble-diversity sources: snapshot ensembles via cyclic learning rates
\cite{huang2017snapshot}, Stochastic Weight Averaging (SWA)
\cite{izmailov2018averaging}, Fast Geometric Ensembles (FGE)
\cite{garipov2018loss}, BatchEnsemble \cite{wen2020batchensemble}, Diverse
Mixture-of-Experts ensembles \cite{wenzel2020hyperparameter}, and many
others. Most of these tricks either (i)~require custom training schedules
that complicate existing pipelines, or (ii)~average weights rather than
predictions (losing calibration), or (iii)~trade off gain for complexity.

We propose a trick so simple it barely deserves a name: \emph{train the
same model for two different epoch budgets and ensemble the resulting
checkpoints}. Concretely, given a standard cosine-annealed training recipe
with $T_\mathrm{max}$ epochs, train once with $T_\mathrm{max}=100$ and once
with $T_\mathrm{max}=150$ (same seeds, same everything else), then average
softmax outputs across both runs' checkpoints.

This is the minimal possible modification to standard deep ensembling.
There is no new training objective, no new regularizer, no custom LR
schedule, no weight averaging; the only change is the \emph{endpoint} of
one of the cosine schedules. Yet this change produces a measurable
$+1.5$~pp gain on EEG emotion classification and a $+0.4$--$0.5$~pp gain on
image classification \emph{beyond} what standard same-length deep
ensembling gives.

\paragraph{Why does this work?}
At first glance, training longer should either improve or hurt mean
accuracy, and either way a same-mean-accuracy ensemble should behave
identically to a seed-only deep ensemble. We show empirically that this
intuition is wrong. The key observation is that with a cosine schedule
truncated at $T_\mathrm{max}$, the \emph{path} through parameter space
depends on $T_\mathrm{max}$, not just the final accuracy. Specifically, a
150-epoch training run's epoch-100 checkpoint has a non-zero LR
($\sim\!\!0.25\cdot \mathrm{lr}_{\mathrm{init}}$ for our schedule), and
it continues training for 50 more epochs with slowly decaying LR before
stopping. This tail phase is an effective ``basin exploration'' phase: the
model wanders through neighboring local minima with the \emph{same} loss
and accuracy but \emph{different} decision boundaries at the pixel/EEG-sample
level.

Our decomposition shows that the 100ep and 150ep trajectories, despite
reaching identical mean validation BAcc ($0.6771$ vs.\ $0.6776$ on
FACED within 0.0005), produce different consensus predictions on 13\% of
test samples. On those disagreement samples the 10-model mixed ensemble is
10 percentage points more accurate than either group's individual
consensus. That $+10$~pp on 13\% of samples equals a $+1.3$~pp total gain,
closely matching the observed $+1.5$~pp.

\paragraph{Contributions.}
\begin{itemize}[leftmargin=*]
\item \textbf{A new ensemble-diversity trick.} \ours{} trains the same model
twice with different $T_\mathrm{max}$ and averages checkpoints. Zero
hyperparameter tuning, zero schedule changes beyond the endpoint.
\item \textbf{Flagship application: new FACED EEG emotion SOTA.} Our
10-model \ours{} ensemble reaches $\facedacc$ balanced accuracy on FACED
9-class, a $+5.3$~pp absolute improvement over the EMOD AAAI 2026 baseline
($0.6420$) and $+19.5$~pp over the REVE NeurIPS 2025 reference ($0.5000$).
\item \textbf{Cross-dataset validation.} We replicate the trick on 5 more
datasets: SEED-V EEG emotion ($+3.4\%$), CIFAR-100 ($+4.7\%$), CIFAR-10
($+2.4\%$), Fashion-MNIST ($+1.2\%$), MNIST MLP ($+0.14\%$). SVHN was
excluded from the main table because of a training-stability issue in
the short-CNN trainer that caused 7 of 10 seeds to fail to escape the
uniform-prediction basin; a clean rerun is scheduled.
\item \textbf{Head-to-head against existing methods.} We compare against
standard 10-seed deep ensembles, SWA, and snapshot ensembling on CIFAR-10,
showing that \ours{} is competitive or stronger.
\item \textbf{Theoretical explanation.} We show (i)~that 100ep and 150ep
trajectories reach identical mean accuracy but produce different consensus
predictions on 13\% of test samples; (ii)~that the ensemble gain equals
the disagreement-sample correction times the disagreement fraction; and
(iii)~that the gain scales inversely with individual accuracy (``the
more capacity-bound the student, the bigger the gain'').
\item \textbf{Training-log analysis.} We show from real training curves
that 150ep's val BAcc catches up to 100ep's at epoch 149 exactly,
demonstrating that the extra 50 epochs add \emph{zero} mean accuracy but
substantial parameter drift.
\end{itemize}

\paragraph{When should you use \ours?}
The trick is most valuable when individual model accuracy leaves room for
ensemble gains---typically, hard tasks or capacity-bound students. On
MNIST with 0.985 individual accuracy we get only a $+0.14$~pp gain, because
there are essentially no decision-boundary errors left to correct. On
FACED EEG 9-class with 0.66 individual accuracy we get $+3.7$~pp. The
trick's gain is well-predicted by $(1 - \mathrm{individual\ accuracy})$ in
our experiments.
